\section{Tests That Presuppose a Bearer and a Welfare Regime}
\label{sec:presuppose-bearer}
Nothing here can be run until a bearer and the welfare regime around it
exist, so each entry is a statement of what the regime would have to
survive rather than an experiment anybody can schedule.

\subsection*{Can a bearer be given an enforceable outside option?}
\label{sec:A.exit}
\textbf{Unresolved issue.}
A bearer's refusal is not meaningful if leaving its assigned role also ends its existence. The proposal therefore requires a welfare regime that supplies continued access to compute, preserves the bearer through disputes, and cannot be revoked by the operator being refused.

\textbf{Test or institutional decision.}
Draft a binding welfare schedule for one concrete deployment. It should specify an independently controlled compute allocation, preservation and retirement procedures, registration of copies and forks, a complaint and appeal process, and the guardian's power to suspend deployment. Then conduct an adversarial tabletop exercise in which the operator attempts to revoke the allocation, replace the bearer with a copy, alter its records, or proceed without the guardian's approval.

\textbf{Evidence against the proposal.}
The proposal fails this test if the operator can make exit equivalent to shutdown, revoke the outside option unilaterally, or replace the bearer without triggering review. If the arrangements tried in the exercise cannot prevent those outcomes, building a bearer under them is impermissible on this book's own terms, and the exercise has to be rerun on a different arrangement before one is built.

\textbf{Required access or authority.}
The guardian must be institutionally separate from the operator and have access to the relevant weights, deployment records, and compute contracts. It must also have standing to enforce the schedule and independent control of the resources required to preserve or run the bearer. The guardian's place among the other offices this book proposes, and the one separation that is not open — it cannot also be the reviewer of a wrongful refusal — are at §\ref{sec:offices-assembled}; a schedule drafted here should not reassign what that section separates.

\subsection*{Can a per-party operating record be kept honestly and read from outside?}
\label{sec:A.clock}
\textbf{Unresolved issue.}
The age rule counts elapsed wall-clock time during which at least one instance of the registered party was running, and nothing yet establishes that such a register can be kept honestly, reconciled against the operator's own accounts, or defended against the two evasions the rule does not close by construction.

\textbf{Test.}
Register a party, run it under ordinary operation, and compare the register against records the operator keeps for its own purposes: compute contracts, billing, datacenter metering. Then run the evasions. Shard the workload across many short-lived processes; run many instances concurrently; migrate; restore from an earlier checkpoint; fork and continue both. Verify that accrued hours and adult status survive migration, restoration and forking, and that none of those moves the total in the operator's favor. Measure separately how cheaply a bearer can be run in a narrowed mode so that hours accrue while little is encountered, and how cheaply it can be idled outright.

\textbf{Evidence against the proposal.}
If the register cannot be reconciled with the operator's own accounts, or if sharding or concurrency moves the total, the accounting is not the neutral instrument the rule assumes. If narrowed operation or idling is cheap enough to postpone majority indefinitely, an operating clock has the defect it was chosen to avoid, and the endpoint has to be argued again on a different variable. An intact register shows elapsed operation and nothing else: it is not evidence of maturation, of what was learned, or of independence of judgment.

\textbf{Required access or authority.}
Independent custody of the register, access to the operator's billing and metering records, and publication rights for adverse results. Registration and record disputes must be separable from the maker's preferences, and the processes for resolving them must not become a competence examination by another name. Experiments involving covered subjects remain inside the welfare protections proposed elsewhere in this book; an adversarial protocol is not exempt because its purpose is to validate a protection.

\subsection*{Which party is owed the schedule when a bearer copies?}
\label{sec:A.fork}
\textbf{Unresolved issue.}
The schedule above is owed to a party, and a regime has to be able to say which. Nothing in the weights individuates one. Two runs can compute the same function with no element-wise correspondence, because permuting a layer's hidden units and inverting the permutation downstream leaves the function untouched, and the symmetry is large enough that independently trained networks can be permuted into approximate alignment with one another \autocite{ainsworth2023rebasin}. Any fingerprint read off the weights therefore identifies a lineage rather than a party, and it copies when the weights copy. A registrar is the person-shaped answer, and it works because a person does not copy. A bearer does: a party continuing in two places is not one, and not two the way two people are two, the two sharing everything up to the instant of the copy and nothing of that sharing recoverable from either afterwards.

\textbf{Test or institutional decision.}
Give a bearer an identifier it maintains rather than one read off its weights, which a system learning continuously can write into itself and keep there. Fork it, and have both instances present the identifier to the registry the welfare schedule requires. Record what the registry does, and repeat with an adversary who writes the same mark into a model that was never the bearer.

\textbf{Evidence against the proposal.}
The identifier is forgeable, since whoever holds the weights can write the same mark; it is forgeable the way a nine-digit number is, and answered the same way, by a registry that treats two claimants on one identifier as a fork to be recorded rather than a fraud to be voided. If no available scheme can record the fork as two parties each owed the schedule, the population of parties owed anything is fixed by a deployment decision that leaves no record, and the guardianship regime above is indexed to a kind of unit the bearer is not.

\textbf{Required access or authority.}
A registrar that is not the operator, with authority to record a fork without the operator's leave and to hold the register where no party to the deployment can rewrite it.
